% ЗГЕНЕРОВАНО scripts/build_format_v2_chapter.py — не редагувати вручну.
% Джерело: editorial/cards/15.*.json, format-v2/rozdil-15.json, paraleli.json.

\naukchapter{15}{20}{Глава 15. Живи легко}
\naukintro{Пʼятнадцята глава — двадцять віршів, і в ній кілька разів звучить те саме, що стало її назвою: живи легко. Це не заклик до безтурботності, а висновок із попереднього: якщо ти не тіло й не діяч (15.4–15.5), то приріст і втрата до тебе не доходять (15.10–15.11). Глава довга, бо повторює ту саму думку з різних боків, і в кінці доходить до найрадикальнішого: відкинути саму медитацію (15.20).}
\vsec{15.1}
\vdeva{यथातथोपदेशेन कृतार्थः सत्त्वबुद्धिमान् । आजीवम् अपि जिज्ञासुः परस् तत्र विमुह्यति ॥}
\vtranslit{ятхататхопадешена крітартхах саттвабуддгіман | аджівам апі джіджнясух парас татра вімухяті ||}
\vperek{Хто має ясний розум, той здійснює мету настановою «як є»; а допитливий, хоч би все життя шукав, — інший у цьому збивається з пуття.}
\vcom{Построкове зіставлення: редакторська чернетка, не рецензовано.}
\vcom{Рядок а}
\vdeva{यथातथोपदेशेन कृतार्थः सत्त्वबुद्धिमान्}
\vtranslit{ятхататхопадешена крітартхах саттвабуддгіман}
\vperek{Хто має ясний розум, той здійснює мету настановою «як є»;}
\vcom{yathātathopadeśena — орудний відмінок «через настанову» з неоднозначним yathā-tathā: у чинному українському перекладі «як є», проте можливе й інше тлумачення способу навчання; kṛtārthaḥ — «той, хто досяг мети», sattvabuddhimān — «той, чий розум ясний». Граматично це описаний стан учня, а не емпіричне твердження, ніби освіта непотрібна.}
\vcom{Рядок б}
\vdeva{आजीवम् अपि जिज्ञासुः परस् तत्र विमुह्यति}
\vtranslit{аджівам апі джіджнясух парас татра вімухяті}
\vperek{а допитливий, хоч би все життя шукав, — інший у цьому збивається з пуття.}
\vcom{ājīvam api — «хоч би протягом життя», jijñāsuḥ — «той, хто прагне знати», paraḥ (у сандгі paras) — «інший», tatra — «тут/у цьому», vimuhyati — «плутається, збивається з пуття». Протиставлення ясності й нескінченного пошуку належить голосу вірша; воно саме собою не засуджує всіх дослідників чи навчання.}
\vcom{`sattvabuddhimān` — «той, чий розум ясний» — стоїть як умова, а не як наслідок: спершу ясність, потім настанова. `yathātathā` («як є») я читаю як прислівник способу при `upadeśa`; інше можливе прочитання — «настанова „так-не-так“», тобто безладний навчальний текст, і тоді сенс інший. Друга половина протиставляє не знання незнанню, а **ясність допитливості**: `jijñāsu` тут не похвала. Це узгоджується з 9.2, де втихає саме `bubhutsā`.}
\vpar{Українська паралель (наша інтерпретація)}{Памфіл Юркевич, «Серце і його значення в духовному житті людини» (1860).}{Наша паралель: коли вірш протиставляє ясний розум допитливому, він ставить настанову вище за пошук. Юркевич теж відводить пізнанню місце, підпорядковане цілому, а не навпаки. Наше порівняння — за будовою думки.}
\vfig{ilustratsii/15-1-motion.pdf}{Дві дороги: настанова як є і довгий пошук}{Схема наша, не історична ілюстрація.}

\vsec{15.2}
\vdeva{मोक्षो विषयवैरस्यं बन्धो वैषयिको रसः । एतावद् एव विज्ञानं यथेच्छसि तथा कुरु ॥}
\vtranslit{мокшо вішаявайрасям бандго вайшаїко расах | етавад ева віджнянам ятхеччасі татха куру ||}
\vperek{Визволення — коли чуттєве втрачає свій смак; неволя — коли воно й далі вабить. Ось і все розуміння. Чини, як бажаєш.}
\vcom{Построкове зіставлення: редакторська чернетка, не рецензовано.}
\vcom{Рядок а}
\vdeva{मोक्षो विषयवैरस्यं बन्धो वैषयिको रसः}
\vtranslit{мокшо вішаявайрасям бандго вайшаїко расах}
\vperek{Визволення — коли чуттєве втрачає свій смак; неволя — коли воно й далі вабить.}
\vcom{mokṣo (mokṣaḥ) — «визволення», viṣayavairasyaṃ — «втрата смаку до чуттєвих предметів»; bandho (bandhaḥ) — «зв'язаність, неволя», vaiṣayiko rasaḥ — «смак до чуттєвого». Рима понять працює через rasa — смак, а не заборону самих матеріальних речей. Це філософська дефініція у творі, а не медичний опис втрати смакових відчуттів.}
\vcom{Рядок б}
\vdeva{एतावद् एव विज्ञानं यथेच्छसि तथा कुरु}
\vtranslit{етавад ева віджнянам ятхеччасі татха куру}
\vperek{Ось і все розуміння. Чини, як бажаєш.}
\vcom{etāvad eva vijñānaṃ — «ось і все розуміння/розрізнення»; yathecchasi (yathā icchasi) — «як бажаєш», tathā kuru — «так чини», наказовий спосіб від kṛ. Формула залишає простір дії після пояснення визволення, але в санскриті немає окремого слова «анархія» чи безумовної політичної програми; будь-яке таке зіставлення було б сучасною інтерпретацією.}
\vcom{Протиставлення viṣaya-vairasyaṃ і vaiṣayika rasaḥ побудовано на образі «смаку»: прив'язаність визначається не самим існуванням предметів, а ставленням до них. Фінальне yathecchasi tathā kuru стоїть після etāvad eva vijñānaṃ; тому «чини, як бажаєш» не слід відривати від попереднього визначення визволення. Формула не містить слова «анархія» і сама собою не обґрунтовує твердження про політичне вчення; порівняння з анархізмом належить до окремого сучасного інтерпретаційного есе, а не до перекладу. Джерельна база цієї чернетки: SARIT verse\_15.2 і GRETIL Avg\_15.2; незалежна санскритологічна рецензія ще потрібна.}
\vpar{Українська паралель (наша інтерпретація)}{Григорій Сковорода, «Убогий Жайворонок».}{Наша паралель: у Сковороди смак до речей теж є мірилом неволі, а не самі речі. Наше порівняння — типологічне, не твердження про джерело.}
\vfig{ilustratsii/15-2-balance.pdf}{Дві шальки: смак чуттєвого і його відсутність}{Схема наша, не історична ілюстрація.}

\vsec{15.3}
\vdeva{वाग्मिप्राज्ञमहोद्योगं जनं मूकजडालसम् । करोति तत्त्वबोधो ऽयम् अतस् त्यक्तो बुभुक्षुभिः ॥}
\vtranslit{ваґміпраджнямаходйоґам джанам мукаджадаласам | кароті таттвабодго ’ям атас тякто бубгукшубгіх ||}
\vperek{Це пізнання суті робить красномовних, тямущих і вельми завзятих людей німими, тупими й лінивими — тому голодні його відкидають.}
\vcom{Найдошкульніше місце глави, і воно **сатиричне**: пізнання суті не підносить людину, а відбирає в неї красномовство, тямущість і завзятість. `bubhukṣubhiḥ` — «голодними»: тими, хто заробляє словом і зусиллям. Тобто текст сам пояснює, чому його вчення не користується попитом, і робить це без жалю. Переклад мусить зберегти різкість трьох пар, не пом'якшуючи їх до «мовчазних, спокійних, бездіяльних».}
\vpar{Українська паралель (ПОТРЕБУЄ ПЕРЕВІРКИ)}{українська барокова проповідь, не визначено.}{ПОТРЕБУЄ ПЕРЕВІРКИ: думка, що знання суті робить красномовних німими, близька до барокової проповіді про марність учення, але конкретного тексту ми не перевірили.}
\vfig{ilustratsii/15-3-neti.pdf}{Красномовні й завзяті стають німими}{Схема наша, не історична ілюстрація.}

\vsec{15.4}
\vdeva{न त्वं देहो न ते देहो भोक्ता कर्ता न वा भवान् । चिद्रूपो ऽसि सदा साक्षी निरपेक्षः सुखं चर ॥}
\vtranslit{на твам дехо на те дехо бгокта карта на ва бгаван | чідрупо ’сі сада сакші нірапекшах сукхам чара ||}
\vperek{Ти не тіло, і тіло не твоє; і не ти — той, хто споживає чи чинить. Ти маєш природу свідомості, повсякчас свідок, без залежності — живи легко.}
\vcom{`sākṣin` («свідок») тут єдиний раз у главі й стоїть у трійці з `cidrūpa` і `nirapekṣa`; переклад мусить тримати його як термін, а не як метафору. `sukhaṃ cara` — наказовий спосіб, «живи легко»: це не опис стану, а запрошення. Формула «не тіло, і тіло не твоє» та сама, що в 8.4, 11.6 і 2.22 — тут до неї додано заперечення **діяча й споживача**, чого в тих місцях не було.}
\vpar{Українська паралель (наша інтерпретація)}{Григорій Сковорода, «Наркіс. Розмова про те: Пізнай себе».}{Наша паралель: у Сковороди справжня людина теж не є ані тілом, ані тим, хто споживає. Наше порівняння — типологічне.}
\vfig{ilustratsii/15-4-neti.pdf}{Тіло й споживач відкинуті: лишається свідок}{Схема наша, не історична ілюстрація.}

\vsec{15.5}
\vdeva{रागद्वेषौ मनोधर्मौ न मनस् ते कदाचन । निर्विकल्पो ऽसि बोधात्मा निर्विकारः सुखं चर ॥}
\vtranslit{раґадвешау манодгармау на манас те кадачана | нірвікалпо ’сі бодгатма нірвікарах сукхам чара ||}
\vperek{Пристрасть і відраза — властивості ума, і ум ніколи не твій. Ти без домислювання, твоя суть — усвідомлення, ти без видозмін — живи легко.}
\vcom{Ланцюг тут двоступеневий: спершу `rāga`/`dveṣa` віднесено до `manas`, а потім заперечено сам `manas` як власність. `nirvikāra` («без видозмін») — те саме слово, що в 11.1; словник спільний. `sukhaṃ cara` повторює 15.4 дослівно, і в перекладі це має бути той самий вислів — інакше зникне наскрізний рефрен глави.}
\vpar{Українська паралель (наша інтерпретація)}{Микола Хвильовий, «Я (Романтика)» (1924).}{Наша паралель: у Хвильового пристрасть і відраза теж показані як властивості ума, а не як голос самої речі. Наше порівняння — за способом бачити.}
\vfig{ilustratsii/15-5-balance.pdf}{Дві шальки: пристрасть і відраза — обидві не твої}{Схема наша, не історична ілюстрація.}

\vsec{15.6}
\vdeva{सर्वभूतेषु चात्मानं सर्वभूतानि चात्मनि । विज्ञाय निरहंकारो निर्ममस् त्वं सुखी भव ॥}
\vtranslit{сарвабгутешу чатманам сарвабгутані чатмані | віджняя нірахамкаро нірмамас твам сукхі бгава ||}
\vperek{Пізнавши себе у всіх істотах і всі істоти — у собі, ти без «я-роблю», без «моє» — будь щасливий.}
\vcom{Двобічність тут обов'язкова: `ātmānam` у всіх істотах **і** всі істоти в `ātman`. Переклад, що зберіг лише один напрямок, втрачає суть. `nirahaṃkāra`/`nirmama` — не «скромний» і не «безкорисливий», а буквально «без `ahaṃkāra`» і «без `mama`»: два займенники, від яких відмовляються. `sukhī bhava` — наказовий спосіб, як і `sukhaṃ cara` у 15.4–15.5.}
\vpar{Українська паралель (наша інтерпретація)}{Митрополит Іларіон Київський, «Слово про закон і благодать» (XI ст.).}{Наша паралель: у київському Слові теж є думка, що людина впізнає себе не в собі одній. Наше порівняння — типологічне; датування цієї книги науково не встановлене.}
\vfig{ilustratsii/15-6-witness.pdf}{Взаємне перебування: ти в усіх, усі в тобі}{Схема наша, не історична ілюстрація.}

\vsec{15.7}
\vdeva{विश्वं स्फुरति यत्रेदं तरङ्गा इव सागरे । तत्त्वम् एव न सन्देहश् चिन्मूर्ते विज्वरो भव ॥}
\vtranslit{вішвам спхураті ятредам таранґа іва саґаре | таттвам ева на сандехаш чінмурте віджваро бгава ||}
\vperek{Те, в чому цей всесвіт постає, як хвилі в морі, — це саме ти, немає сумніву, о втілення свідомості; будь без трепету.}
\vcom{Образ хвиль і моря — єдиний такий у книзі, і він робить роботу, яку не робить жодне означення: хвиля не **відмінна** від моря й не **тотожна** йому як річ. `cinmūrte` — звертання «о втілення свідомості», споріднене з `cidrūpa` у 15.4; переклад мусить тримати спільний корінь. `vijvara` («без трепету») — від `jvara`, гарячка; це не спокій, а відсутність пропасниці.}
\vpar{Українська паралель (наша інтерпретація)}{Богдан-Ігор Антонич, «Ротації» (1938).}{Наша паралель: в Антонича всесвіт теж постає як хвиля в морі, а не як споруда. Наше порівняння — за образом.}
\vfig{ilustratsii/15-7-motion.pdf}{Всесвіт як хвилі в морі: це саме ти}{Схема наша, не історична ілюстрація.}

\vsec{15.8}
\vdeva{श्रद्धत्स्व तात श्रद्धत्स्व नात्र मोहं कुरुष्व भोः । ज्ञानस्वरूपो भगवान् आत्मा त्वं प्रकृतेः परः ॥}
\vtranslit{шраддгатсва тата шраддгатсва натра мохам курушва бгох | джнянасварупо бгаґаван атма твам пракрітех парах ||}
\vperek{Повір, любий, повір; не роби тут затьмарення, друже. Твоя суть — знання; бгаґаван ātman — ти, понад природою.}
\vcom{Єдине місце в книзі, де вчення вимагає **віри** — і вимагає двічі, повтором `śraddhatsva`. Це суперечить загальній настанові тексту «перевір сам», і я не згладжую цю суперечність: вона є. `bhagavān ātmā` — рідкісне поєднання, де `ātman` названо `bhagavant`; переклад мусить зберегти обидва слова, не зводячи до «Бог у тобі». `prakṛteḥ paraḥ` — «понад природою», тобто не складник її.}
\vpar{Українська паралель (ПОТРЕБУЄ ПЕРЕВІРКИ)}{українська релігійна поезія, не визначено.}{ПОТРЕБУЄ ПЕРЕВІРКИ: пряме звертання повір, любий із повторенням близьке до української релігійної поезії, але конкретного твору ми не звірили.}
\vfig{ilustratsii/15-8-flame.pdf}{Полумʼя знання: повір — і не роби затьмарення}{Схема наша, не історична ілюстрація.}

\vsec{15.9}
\vdeva{गुणैः संवेष्टितो देहस् तिष्ठत्य् आयाति याति च । आत्मा न गन्ता नागन्ता किम् एनम् अनुशोचसि ॥}
\vtranslit{ґунайх самвештіто дехас тіштхатй аяті яті ча | атма на ґанта наґанта кім енам анушочасі ||}
\vperek{Тіло, обгорнуте гунами, стоїть, приходить і йде; ātman не йде і не приходить — чого ти за ним жалкуєш?}
\vcom{Заперечення подвійне й симетричне: `na gantā na agantā` — не йде **і не приходить**, тобто не підпадає під жоден бік пари. Питання `kim enam anuśocasi` — пряме звертання до того, хто сумує; у главі це вже друге таке звертання після 15.6, і разом вони задають тон: не пояснення, а докір. `anuśocasi` — «жалкуєш за ним», не «сумуєш» взагалі.}
\vpar{Українська паралель (наша інтерпретація)}{Тарас Шевченко, «Марія» (1859).}{Наша паралель: у Шевченка те, що не йде й не приходить, теж протиставлене тілу, яке приходить і відходить. Наше порівняння — типологічне, не про джерело.}
\vfig{ilustratsii/15-9-motion.pdf}{Тіло приходить і йде, а те, що бачить, стоїть}{Схема наша, не історична ілюстрація.}
\vfig{ilustratsii-naratyvni/15-9-tilo-ide.pdf}{Тіло йде дорогою, лишаючи сліди, а той, хто сидить, не йде}{Наша ілюстрація до вірша 15.9 — не історичне зображення.}

\vsec{15.10}
\vdeva{देहस् तिष्ठतु कल्पान्तं गच्छत्य् अद्यैव वा पुनः । क्व वृद्धिः क्व च वा हानिस् तव चिन्मात्ररूपिणः ॥}
\vtranslit{дехас тіштхату калпантам ґаччатй адяйва ва пунах | ква вріддгіх ква ча ва ханіс тава чінматрарупінах ||}
\vperek{Нехай тіло стоїть до кінця світового періоду — або нехай іде хоч сьогодні. Де приріст і де втрата для тебе, чия природа — сама лише свідомість?}
\vcom{Пара наказових способів (`tiṣṭhatu` … `gacchatu`) задає байдужість до тривалості життя: обидва варіанти рівноцінні. `kalpānta` («кінець світового періоду») — не поетична гіпербола, а максимально можливий строк, і саме тому друга частина звучить різко: `adya eva` («хоч сьогодні»). `cinmātrarūpa` («чия природа — сама лише свідомість») — те саме, що `cinmūrte` у 15.7 і `cidrūpa` у 15.4; словник лишається спільним.}
\vpar{Українська паралель (наша інтерпретація)}{Ліна Костенко, «Сад нетанучих скульптур» (1987).}{Наша паралель: у Костенко питання про приріст і втрату теж виявляється поставленим не до тієї міри. Наше порівняння — за будовою думки.}
\vfig{ilustratsii/15-10-balance.pdf}{Дві шальки: тіло надовго і тіло на день}{Схема наша, не історична ілюстрація.}

\vsec{15.11}
\vdeva{त्वय्य् अनन्तमहाम्भोधौ विश्ववीचिः स्वभावतः । उदेतु वास्तम् आयातु न ते वृद्धिर् न वा क्षतिः ॥}
\vtranslit{твайй анантамахамбгодгау вішвавічіх свабгаватах | удету вастам аяту на те вріддгір на ва кшатіх ||}
\vperek{У тобі, безмежному великому океані, хвиля всесвіту за власною природою нехай піднімається або спадає — немає в тебе ні приросту, ні втрати.}
\vcom{Той самий образ, що в 15.7, але розгорнутий: там хвилі були в морі, тут ти **сам** — океан, і всесвіт — одна хвиля в ньому. `svabhāvataḥ` («за власною природою») важливе: хвиля не робиться нею ззовні, вона така за природою. Пара `vṛddhi`/`kṣati` повторює 15.10 дослівно — це рефрен двох віршів, і в перекладі він мусить звучати однаково.}
\vpar{Українська паралель (наша інтерпретація)}{Олександр Довженко, «Поема про море» (1956).}{Наша паралель: у Довженка хвиля теж не додає й не забирає нічого в моря, у якому вона піднялася. Наше порівняння — за образом.}
\vfig{ilustratsii/15-11-motion.pdf}{Хвиля всесвіту в безмежному океані}{Схема наша, не історична ілюстрація.}

\vsec{15.12}
\vdeva{तात चिन्मात्ररूपो ऽसि न ते भिन्नम् इदं जगत् । अतः कस्य कथं कुत्र हेयोपादेयकल्पना ॥}
\vtranslit{тата чінматрарупо ’сі на те бгіннам ідам джаґат | атах кася катхам кутра хейопадеякалпана ||}
\vperek{Любий, твоя природа — сама лише свідомість; цей світ не відмінний від тебе. Тому — для кого, як і де вигадувати те, що відкидати й приймати?}
\vcom{Три питання підряд (`kasya katham kutra`) — не риторичний повтор, а розбір по трьох осях: для кого, яким чином, де. `kalpanā` — не «вибір», а **вигадування**: текст не каже, що вибір між добрим і поганим складний, він каже, що його нема кому робити. `tāta` — те саме ласкаве звертання, що в 9.2 і 15.8.}
\vpar{Українська паралель (наша інтерпретація)}{Іван Вишенський, «Книжка» (кінець XVI ст.).}{Наша паралель: у Вишенського світ теж не відділений від того, хто його бачить, і саме тому вибір між прийняти й відкинути втрачає предмет. Наше порівняння — типологічне.}
\vfig{ilustratsii/15-12-neti.pdf}{Відкидати й приймати нічого: світ не відмінний від тебе}{Схема наша, не історична ілюстрація.}

\vsec{15.13}
\vdeva{एकस्मिन् अव्यये शान्ते चिदाकाशे ऽमले त्वयि । कुतो जन्म कुतो कर्म कुतो ऽहंकार एव च ॥}
\vtranslit{екасмін ав'яє шанте чідакаше ’мале тваї | куто джанма куто карма куто ’хамкара ева ча ||}
\vperek{У тобі — одному, незмінному, втихлому, незаплямованому просторі свідомості — звідки народження, звідки діяння, звідки ж «я-роблю»?}
\vcom{`cidākāśa` («простір свідомості») — єдине таке складне слово в книзі; переклад мусить лишити його як термін, а не як метафору. Чотири означення (`eka`, `avyaya`, `śānta`, `amala`) звужують значення до одного, і кожне наступне сильніше. Три `kutaḥ` ідуть у порядку виникнення: народження → діяння → `ahaṃkāra`, і останнє поставлене як те, що робить перші два можливими.}
\vpar{Українська паралель (ПОТРЕБУЄ ПЕРЕВІРКИ)}{українська містична проза, не визначено.}{ПОТРЕБУЄ ПЕРЕВІРКИ: питання, звідки народження й діяння в тому, що незмінне, найближче до української містичної прози, але назвати твір без звірки ми не можемо.}
\vfig{ilustratsii/15-13-space.pdf}{Порожній простір свідомості: звідки народження й діяння?}{Схема наша, не історична ілюстрація.}

\vsec{15.14}
\vdeva{यद् त्वं पश्यसि तत्रैकस् त्वम् एव प्रतिभाससे । किं पृथग् भासते स्वर्णात् कटकाङ्गदनूपुरम् ॥}
\vtranslit{яд твам пашясі татрайкас твам ева пратібгасасе | кім прітхаґ бгасате сварнат катаканґаданупурам ||}
\vperek{Те, що ти бачиш, — там сяєш один саме ти. Чи сяє окремо від золота браслет, наруч або ножна прикраса?}
\vcom{Риторичне питання з відповіддю «ні»: три прикраси названі поіменно, щоб показати, що різниця між ними — лише форма, а речовина одна. Це той самий хід, що з хвилями в 15.7 і 15.11, але тут він зроблений на **речі**, а не на стихії. `pratibhāsase` — «сяєш, постаєш», споріднене з `sphurati` в 15.7; переклад мусить тримати цю спорідненість.}
\vpar{Українська паралель (наша інтерпретація)}{Григорій Сковорода, «Ікона Алківіадська».}{Наша паралель: у Сковороди прикраса й золото теж співвідносяться як видиме й те, що в ньому сяє. Наше порівняння — за образом.}
\vfig{ilustratsii/15-14-thread.pdf}{Прикраси на нитці: усі вони — золото}{Схема наша, не історична ілюстрація.}
\vfig{ilustratsii-naratyvni/15-14-brastelet-i-zoloto.pdf}{Злиток золота на плиті й браслет поруч: окремо від золота браслет не сяє}{Наша ілюстрація до вірша 15.14 — не історичне зображення.}

\vsec{15.15}
\vdeva{अयं सो ऽहम् अयं नाहम् विभागम् इति संत्यज । सर्वम् आत्मेति निश्चित्य निःसंकल्पः सुखी भव ॥}
\vtranslit{аям со ’хам аям нахам вібгаґам іті самтяджа | сарвам атметі нішчітя ніхсамкалпах сукхі бгава ||}
\vperek{«Це — він, це я; це — не я» — відкинь такий поділ. Вирішивши, що все — ātman, будь без побудов і щасливий.}
\vcom{Тут відкидається не одна зі сторін поділу, а **сама навичка ділити**: `vibhāgam … saṃtyaja`. `niḥsaṃkalpa` — не «без думок», а без `saṃkalpa`, наміру-побудови, тобто без попереднього складання того, що має бути. Це те саме слово, що `kalpanā` в 15.12, і переклад мусить бути спорідненим.}
\vpar{Українська паралель (наша інтерпретація)}{Микола Зеров, «Камена» (1924).}{Наша паралель: у Зерова поділ на це і не це теж знімається ясністю, а не зусиллям. Наше порівняння — типологічне.}
\vfig{ilustratsii/15-15-neti.pdf}{Поділ на це і не це відкинуто: усе — ātman}{Схема наша, не історична ілюстрація.}

\vsec{15.16}
\vdeva{तवैवाज्ञानतो विश्वं त्वम् एकः परमार्थतः । त्वत्तो ऽअन्यो नास्ति संसारी नासंसारी च कश्चन ॥}
\vtranslit{тавайваджнянато вішвам твам еках парамартхатах | тватто ’анйо насті самсарі насамсарі ча кашчана ||}
\vperek{Всесвіт — саме через твоє незнання; у справжньому сенсі ти один. Немає іншого, ніж ти — ні того, хто в сансарі, ні того, хто поза нею.}
\vcom{Заперечення знову симетричне: немає ні `saṃsārī`, ні `asaṃsārī` — тобто не існує ні того, хто в сансарі, ні того, хто з неї вийшов. Це найрадикальніше з тверджень глави, і його легко зіпсувати перекладом «немає ні сансари, ні визволення»: у тексті йдеться про **осіб**, а не про стани. `paramārthataḥ` («у справжньому сенсі») — усталений термін, не «насправді» в побутовому значенні.}
\vpar{Українська паралель (ПОТРЕБУЄ ПЕРЕВІРКИ)}{українська філософська проза, не визначено.}{ПОТРЕБУЄ ПЕРЕВІРКИ: думка, що всесвіт постає через незнання, найближча до української філософської прози, але конкретного твору ми не перевірили.}
\vfig{ilustratsii/15-16-space.pdf}{Всесвіт як горщик, а простір — один}{Схема наша, не історична ілюстрація.}

\vsec{15.17}
\vdeva{भ्रान्तिमात्रम् इदं विश्वं न किञ्चिद् इति निश्चयी । निर्वासनः स्फूर्तिमात्रो न किञ्चिद् इव शाम्यति ॥}
\vtranslit{бгрантіматрам ідам вішвам на кіньчід іті нішчаї | нірвасанах спхуртіматро на кіньчід іва шам'яті ||}
\vperek{«Цей всесвіт — саме лише блукання, аніщо» — хто має таке переконання, той без усталених схильностей, саме лише сяйнування, втихає, наче аніщо.}
\vcom{**Цей вірш майже дослівно збігається з 11.8.** Різниця лише в першому складнику: тут `bhrāntimātram` («саме лише блукання»), там `nānāścaryam` («різноманітно дивовижне»). Друга половина — `nirvāsanaḥ sphūrtimātro na kiñcid iva śāmyati` — ідентична в обох місцях. Це не помилка картки й не недогляд: текст справді повторює формулу, змінюючи лише означення світу. Я зберіг переклад другої половини **дослівно той самий**, що в 11.8, щоб повтор було видно. `bhrānti` — «блукання» в сенсі омани-ходіння, не «ілюзія» як технічний термін.}
\vpar{Українська паралель (наша інтерпретація)}{Іван Франко, «Мойсей» (1905).}{Наша паралель: у Франка блукання теж може виявитися тим, що не має власної суті. Наше порівняння — типологічне.}
\vfig{ilustratsii/15-17-illusion.pdf}{Всесвіт під серпанком: блукання, аніщо}{Схема наша, не історична ілюстрація.}

\vsec{15.18}
\vdeva{एक एव भवाम्भोधाव् आसीद् अस्ति भविष्यति । न ते बन्धो ऽस्ति मोक्षो वा कृतकृत्यः सुखं चर ॥}
\vtranslit{ека ева бгавамбгодгав асід асті бгавішяті | на те бандго ’сті мокшо ва крітакрітях сукхам чара ||}
\vperek{Один-єдиний в океані буття — був, є, буде. Немає в тебе ні пут, ні визволення; той, чию справу зроблено, — живи легко.}
\vcom{`na te bandho 'sti mokṣo vā` — та сама симетрична пара, що в 15.16: немає ні пут, ні визволення. Переклад мусить тримати їх **як пару**, бо звичне «ти вільний» знищило б думку. Три часи (`āsīt asti bhaviṣyati`) стоять у ряд, як у 19.3, але там вони скасовувалися, а тут стверджуються. `kṛtakṛtya` — те саме, що `kṛtārtha` в 12.8, і переклад має бути спорідненим.}
\vpar{Українська паралель (наша інтерпретація)}{Василь Стус, «Час творчості».}{Наша паралель: у Стуса той, чию справу зроблено, теж не має ні пут, ні визволення — він просто живе. Наше порівняння — за будовою думки.}
\vfig{ilustratsii/15-18-space.pdf}{Один в океані буття: був, є, буде}{Схема наша, не історична ілюстрація.}

\vsec{15.19}
\vdeva{मा संकल्पविकल्पाभ्यां चित्तं क्षोभय चिन्मय । उपशाम्य सुखं तिष्ठ स्वात्मन्य् आनन्दविग्रहे ॥}
\vtranslit{ма самкалпавікалпабгям чіттам кшобгая чінмая | упашам'я сукхам тіштха сватманй анандавіґрахе ||}
\vperek{Не збурюй ум наміром і домислюванням, о сповнений свідомості. Втихнувши, стій легко у собі — втіленні блаженства.}
\vcom{`saṃkalpa` і `vikalpa` названі парою як те, що збурює: перше — складання наміру, друге — роздвоєння-домислювання. Це ті самі два корені, що в 15.12 (`kalpanā`) і 15.15 (`niḥsaṃkalpa`), і разом вони утворюють послідовний словник глави. `ānandavigraha` («втілення блаженства») — єдине таке поєднання в книзі; це не «блаженна людина», а означення **себе**.}
\vpar{Українська паралель (наша інтерпретація)}{Павло Тичина, «Сонячні кларнети» (1918).}{Наша паралель: у Тичини втихнути й стояти легко теж означає перестати збурювати те, що й так сяє. Наше порівняння — за образом.}
\vfig{ilustratsii/15-19-flame.pdf}{Полумʼя без вітру: втихнувши, стій легко}{Схема наша, не історична ілюстрація.}

\vsec{15.20}
\vdeva{त्यजैव ध्यानं सर्वत्र मा किञ्चिद् धृदि धारय । आत्मा त्वं मुक्त एवासि किं विमृश्य करिष्यसि ॥}
\vtranslit{тяджайва дгянам сарватра ма кіньчід дгріді дгарая | атма твам мукта евасі кім вімрішя карішясі ||}
\vperek{Відкинь саму медитацію, скрізь не тримай нічого в серці. Ти — ātman, ти вже вільний: що ти робитимеш, далі розмірковуючи?}
\vcom{Останній вірш глави відкидає **медитацію** — те, з чого починалася майже вся релігійна практика. `tyaja eva` з підсиленням: не «зменш», а «відкинь саме». `mukta eva asi` — «ти вже вільний», і `eva` тут не прикраса: без нього фраза звучала б як обіцянка. Завершальне питання `kiṃ … kariṣyasi` не закликає до дії — воно показує, що робити нічого. Це узгоджується з 15.3, де пізнання суті відбирає завзятість.}
\vpar{Українська паралель (наша інтерпретація)}{Леся Українка, «Оргія» (1913).}{Наша паралель: у Лесі Українки відкинути саму медитацію — це теж крок, а не втрата, бо мета вже не попереду. Наше порівняння — типологічне.}
\vfig{ilustratsii/15-20-neti.pdf}{Медитація відкинута: ти вже вільний}{Схема наша, не історична ілюстрація.}

